% Options for packages loaded elsewhere
\PassOptionsToPackage{unicode}{hyperref}
\PassOptionsToPackage{hyphens}{url}
\documentclass[
  ignorenonframetext,
  aspectratio=169]{beamer}
\newif\ifbibliography
\usepackage{pgfpages}
\setbeamertemplate{caption}[numbered]
\setbeamertemplate{caption label separator}{: }
\setbeamercolor{caption name}{fg=normal text.fg}
\beamertemplatenavigationsymbolsempty
% remove section numbering
\setbeamertemplate{part page}{
  \centering
  \begin{beamercolorbox}[sep=16pt,center]{part title}
    \usebeamerfont{part title}\insertpart\par
  \end{beamercolorbox}
}
\setbeamertemplate{section page}{
  \centering
  \begin{beamercolorbox}[sep=12pt,center]{section title}
    \usebeamerfont{section title}\insertsection\par
  \end{beamercolorbox}
}
\setbeamertemplate{subsection page}{
  \centering
  \begin{beamercolorbox}[sep=8pt,center]{subsection title}
    \usebeamerfont{subsection title}\insertsubsection\par
  \end{beamercolorbox}
}
% Prevent slide breaks in the middle of a paragraph
\widowpenalties 1 10000
\raggedbottom
\AtBeginPart{
  \frame{\partpage}
}
\AtBeginSection{
  \ifbibliography
  \else
    \frame{\sectionpage}
  \fi
}
\AtBeginSubsection{
  \frame{\subsectionpage}
}
\usepackage{iftex}
\ifPDFTeX
  \usepackage[T1]{fontenc}
  \usepackage[utf8]{inputenc}
  \usepackage{textcomp} % provide euro and other symbols
\else % if luatex or xetex
  \usepackage{unicode-math} % this also loads fontspec
  \defaultfontfeatures{Scale=MatchLowercase}
  \defaultfontfeatures[\rmfamily]{Ligatures=TeX,Scale=1}
\fi
\usepackage{lmodern}
\ifPDFTeX\else
  % xetex/luatex font selection
\fi
% Use upquote if available, for straight quotes in verbatim environments
\IfFileExists{upquote.sty}{\usepackage{upquote}}{}
\IfFileExists{microtype.sty}{% use microtype if available
  \usepackage[]{microtype}
  \UseMicrotypeSet[protrusion]{basicmath} % disable protrusion for tt fonts
}{}
\makeatletter
\@ifundefined{KOMAClassName}{% if non-KOMA class
  \IfFileExists{parskip.sty}{%
    \usepackage{parskip}
  }{% else
    \setlength{\parindent}{0pt}
    \setlength{\parskip}{6pt plus 2pt minus 1pt}}
}{% if KOMA class
  \KOMAoptions{parskip=half}}
\makeatother
\usepackage{color}
\usepackage{fancyvrb}
\newcommand{\VerbBar}{|}
\newcommand{\VERB}{\Verb[commandchars=\\\{\}]}
\DefineVerbatimEnvironment{Highlighting}{Verbatim}{commandchars=\\\{\}}
% Add ',fontsize=\small' for more characters per line
\usepackage{framed}
\definecolor{shadecolor}{RGB}{248,248,248}
\newenvironment{Shaded}{\begin{snugshade}}{\end{snugshade}}
\newcommand{\AlertTok}[1]{\textcolor[rgb]{0.94,0.16,0.16}{#1}}
\newcommand{\AnnotationTok}[1]{\textcolor[rgb]{0.56,0.35,0.01}{\textbf{\textit{#1}}}}
\newcommand{\AttributeTok}[1]{\textcolor[rgb]{0.13,0.29,0.53}{#1}}
\newcommand{\BaseNTok}[1]{\textcolor[rgb]{0.00,0.00,0.81}{#1}}
\newcommand{\BuiltInTok}[1]{#1}
\newcommand{\CharTok}[1]{\textcolor[rgb]{0.31,0.60,0.02}{#1}}
\newcommand{\CommentTok}[1]{\textcolor[rgb]{0.56,0.35,0.01}{\textit{#1}}}
\newcommand{\CommentVarTok}[1]{\textcolor[rgb]{0.56,0.35,0.01}{\textbf{\textit{#1}}}}
\newcommand{\ConstantTok}[1]{\textcolor[rgb]{0.56,0.35,0.01}{#1}}
\newcommand{\ControlFlowTok}[1]{\textcolor[rgb]{0.13,0.29,0.53}{\textbf{#1}}}
\newcommand{\DataTypeTok}[1]{\textcolor[rgb]{0.13,0.29,0.53}{#1}}
\newcommand{\DecValTok}[1]{\textcolor[rgb]{0.00,0.00,0.81}{#1}}
\newcommand{\DocumentationTok}[1]{\textcolor[rgb]{0.56,0.35,0.01}{\textbf{\textit{#1}}}}
\newcommand{\ErrorTok}[1]{\textcolor[rgb]{0.64,0.00,0.00}{\textbf{#1}}}
\newcommand{\ExtensionTok}[1]{#1}
\newcommand{\FloatTok}[1]{\textcolor[rgb]{0.00,0.00,0.81}{#1}}
\newcommand{\FunctionTok}[1]{\textcolor[rgb]{0.13,0.29,0.53}{\textbf{#1}}}
\newcommand{\ImportTok}[1]{#1}
\newcommand{\InformationTok}[1]{\textcolor[rgb]{0.56,0.35,0.01}{\textbf{\textit{#1}}}}
\newcommand{\KeywordTok}[1]{\textcolor[rgb]{0.13,0.29,0.53}{\textbf{#1}}}
\newcommand{\NormalTok}[1]{#1}
\newcommand{\OperatorTok}[1]{\textcolor[rgb]{0.81,0.36,0.00}{\textbf{#1}}}
\newcommand{\OtherTok}[1]{\textcolor[rgb]{0.56,0.35,0.01}{#1}}
\newcommand{\PreprocessorTok}[1]{\textcolor[rgb]{0.56,0.35,0.01}{\textit{#1}}}
\newcommand{\RegionMarkerTok}[1]{#1}
\newcommand{\SpecialCharTok}[1]{\textcolor[rgb]{0.81,0.36,0.00}{\textbf{#1}}}
\newcommand{\SpecialStringTok}[1]{\textcolor[rgb]{0.31,0.60,0.02}{#1}}
\newcommand{\StringTok}[1]{\textcolor[rgb]{0.31,0.60,0.02}{#1}}
\newcommand{\VariableTok}[1]{\textcolor[rgb]{0.00,0.00,0.00}{#1}}
\newcommand{\VerbatimStringTok}[1]{\textcolor[rgb]{0.31,0.60,0.02}{#1}}
\newcommand{\WarningTok}[1]{\textcolor[rgb]{0.56,0.35,0.01}{\textbf{\textit{#1}}}}
\usepackage{longtable,booktabs,array}
\usepackage{calc} % for calculating minipage widths
\usepackage{caption}
% Make caption package work with longtable
\makeatletter
\def\fnum@table{\tablename~\thetable}
\makeatother
\setlength{\emergencystretch}{3em} % prevent overfull lines
\providecommand{\tightlist}{%
  \setlength{\itemsep}{0pt}\setlength{\parskip}{0pt}}
\usepackage{amsmath}
\usepackage{xcolor}
\setbeamercolor{frametitle}{fg=black}
\usepackage{graphicx}
\usebackgroundtemplate{\includegraphics[width=\paperwidth]{rmdFigs/RH_template_Page_2.png}}
\addtobeamertemplate{frametitle}{\vspace*{.25in}}{\vspace*{.25in}}
\setbeamerfont{frametitle}{size=\huge}
\usepackage{tikz}
\usepackage{bookmark}
\IfFileExists{xurl.sty}{\usepackage{xurl}}{} % add URL line breaks if available
\urlstyle{same}
\hypersetup{
  pdftitle={Communicating Science using R Markown},
  hidelinks,
  pdfcreator={LaTeX via pandoc}}

\title{Communicating Science using R Markown}
\author{W. Evan Johnson, Ph.D.\\
Professor, Division of Infectious Disease\\
Director, Center for Data Science\\
Rutgers University -- New Jersey Medical School\\
\href{mailto:w.evan.johnson@rutgers.edu}{\nolinkurl{w.evan.johnson@rutgers.edu}}}
\date{2026-02-09}

\begin{document}
\frame{\titlepage}

\begin{frame}{Communicating Science}
\phantomsection\label{communicating-science}
Effective science communicators educate non-specialist audiences about
scientific topics, issues, and debates in ways that are informative,
accessible, and empowering.

Science communicators should be able to answer the following questions:

\begin{itemize}
\tightlist
\item
  Who is my audience?
\item
  What is my message for my audience?
\item
  What medium am I going to use to communicate my message to my
  audience?
\end{itemize}
\end{frame}

\begin{frame}{Communicating Science}
\phantomsection\label{communicating-science-1}
The \textbf{lay public} is made up of all the people who are not experts
in a specific field. When addressing the lay public:

\begin{itemize}
\tightlist
\item
  Keep the story simple and front-load exciting aspects
\item
  Make the story relevant to your audience
\item
  Use analogies and visuals
\item
  Front-load the story
\item
  Avoid jargon - simple language but don't oversimplify
\item
  Include the people and the process (challenges, successes,
  collaborations, etc.)
\end{itemize}
\end{frame}

\begin{frame}{Communicating Science}
\phantomsection\label{communicating-science-2}
The \textbf{media} is a ``mediator'' between scientists and the public.
Note the media is not a homogenous group.

Describing your process, challenges, successes, and collaborations are
important for writing an informative and engaging press release. Read a
few popular science articles to get a sense of how your research might
eventually appear in the news and magazines.

Articles about your work should include visuals--videos or photos--that
will draw readers' attention to the article and help them grasp the gist
of the piece.
\end{frame}

\begin{frame}{Communicating Science}
\phantomsection\label{communicating-science-3}
Scientists can share their knowledge with \textbf{policy makers} through
meetings, testimonies, and open presentations. Suggestions for
communicating with policy makers:

\begin{itemize}
\tightlist
\item
  Know what issues policy makers are currently discussing and debating
\item
  Keep your explanations simple and relevant
\item
  Think of some actionable solutions to the problem
\item
  Think about the problem and solution in the context of the policy
  maker's constituency
\item
  Be confident in yourself and what you know
\item
  Approach a meeting as a conversation, not a presentation
\item
  Create a one-pager with your message and key points
\end{itemize}
\end{frame}

\begin{frame}{Communicating Science}
\phantomsection\label{communicating-science-4}
Visuals make the data supporting your message clear and accessible to
your audience. Science visualizations include:

\begin{enumerate}
\tightlist
\item
  Graphs, tables, and infographics
\item
  Conceptual diagrams
\item
  Maps, Satellite photos
\item
  Photographs
\end{enumerate}

Keep the following points in mind:

\begin{enumerate}
\tightlist
\item
  Use a consistent style and format
\item
  Use colors with purpose
\item
  Use high-resolution graphics
\item
  Format your graphics and include labels, legends, and captions
\end{enumerate}
\end{frame}

\begin{frame}{Reproducible Reports (R markdown)}
\phantomsection\label{reproducible-reports-r-markdown}
The final product of a data analysis project is often a report:
scientific publications, news articles, an analysis report for your
company, or lecture notes for a class.

Now imagine after you are done you realize you:

\begin{itemize}
\tightlist
\item
  had the wrong dataset
\item
  have a new dataset for the same analysis
\item
  made a mistake and fix the error, or
\item
  your boss or someone you are training wants to see the code and be
  able to reproduce the results \vskip .1in
\end{itemize}

Situations like these are common for a data scientist.
\end{frame}

\begin{frame}{R markdown}
\phantomsection\label{r-markdown}
R markdown is a format for \textbf{literate programming} documents.
Literate programming weaves instructions, documentation, equations, and
detailed comments among executable code.

It is based on \textbf{markdown}, a markup language that is widely used
to generate html pages. You can learn more about markdown with the
following tutorial:
\href{https://www.markdowntutorial.com/}{\textbf{click here}}
\end{frame}

\begin{frame}[fragile]{R markdown}
\phantomsection\label{r-markdown-1}
You can start an R markdown document in RStudio by clicking on
\textbf{File}, \textbf{New File}, then \textbf{R markdown}. You will
then be asked for a title and author.

Once you gain experience with R markdown, you will be able to do this
without the template and can simply start from a blank template.

As a convention, we use the \texttt{.Rmd} suffix for R markdown files
\end{frame}

\begin{frame}[fragile]{Compiling the document using \texttt{knitR}}
\phantomsection\label{compiling-the-document-using-knitr}
With R markdown, you need to \textbf{compile} the document into the
final report. We use the \texttt{knitR} package to to do this. The
specific function used to compile is the \texttt{knit} function, which
takes a filename as input.

RStudio provides a \texttt{Knit} button that makes it easy and
convenient to compile the document.
\end{frame}

\begin{frame}[fragile]{The Header}
\phantomsection\label{the-header}
At the top of the document is the R markdown header:

\begin{Shaded}
\begin{Highlighting}[]
\SpecialCharTok{{-}{-}{-}}
\NormalTok{title}\SpecialCharTok{:} \StringTok{"Nanostring Analysis"}
\NormalTok{author}\SpecialCharTok{:} \StringTok{"Evan Johnson"}
\NormalTok{date}\SpecialCharTok{:} \StringTok{"12/5/2019"}
\NormalTok{output}\SpecialCharTok{:}\NormalTok{ html\_document}\SpecialCharTok{:}
\SpecialCharTok{{-}{-}{-}}
\end{Highlighting}
\end{Shaded}

R markdown reports can be to be in HTML, PDF, Microsoft Word, or
presentation formats. By changing the \texttt{output} to, for example
\texttt{pdf\_document} or \texttt{word\_document}, we can control the
type of output that is produced.
\end{frame}

\begin{frame}[fragile]{The Header}
\phantomsection\label{the-header-1}
Other output options include code folding, themes, etc.:

\begin{Shaded}
\begin{Highlighting}[]
\SpecialCharTok{{-}{-}{-}}
\NormalTok{title}\SpecialCharTok{:} \StringTok{"Nanostring Analysis"}
\NormalTok{author}\SpecialCharTok{:} \StringTok{"Evan Johnson"}
\NormalTok{date}\SpecialCharTok{:} \StringTok{"12/5/2019"}
\NormalTok{output}\SpecialCharTok{:}
\NormalTok{  html\_document}\SpecialCharTok{:}
\NormalTok{    code\_folding}\SpecialCharTok{:}\NormalTok{ hide}
\NormalTok{    toc}\SpecialCharTok{:}\NormalTok{ true}
\NormalTok{    toc\_float}\SpecialCharTok{:}\NormalTok{ true}
\NormalTok{    theme}\SpecialCharTok{:} \StringTok{"flatly"}
\NormalTok{editor\_options}\SpecialCharTok{:} 
\NormalTok{  chunk\_output\_type}\SpecialCharTok{:}\NormalTok{ console}
\SpecialCharTok{{-}{-}{-}}
\end{Highlighting}
\end{Shaded}
\end{frame}

\begin{frame}[fragile]{R Code Chunks}
\phantomsection\label{r-code-chunks}
In various places in the document, we see something like this:

\begin{Shaded}
\begin{Highlighting}[]
\NormalTok{\textasciigrave{}\textasciigrave{}\textasciigrave{}\{r\}}
\NormalTok{summary(pressure)}
\NormalTok{\textasciigrave{}\textasciigrave{}\textasciigrave{}}
\end{Highlighting}
\end{Shaded}

These are the \textbf{code chunks}. When you compile the document, the R
code inside the chunk, in this case \texttt{summary(pressure)}, will be
evaluated and the result included in that position in the final
document.

To add your own R chunks, you can type the characters above quickly with
the key binding command-option-I on the Mac and Ctrl-Alt-I on Windows.
\end{frame}

\begin{frame}[fragile]{R Code Chunks}
\phantomsection\label{r-code-chunks-1}
Its a good habit to adding a label to the R code chunks. This will be
very useful when debugging, among other situations. You do this by
adding a descriptive word like this:

\begin{Shaded}
\begin{Highlighting}[]
\NormalTok{\textasciigrave{}\textasciigrave{}\textasciigrave{}\{r pressure{-}summary\}}
\NormalTok{summary(pressure)}
\NormalTok{\textasciigrave{}\textasciigrave{}\textasciigrave{}}
\end{Highlighting}
\end{Shaded}
\end{frame}

\begin{frame}[fragile]{R Code Chunks}
\phantomsection\label{r-code-chunks-2}
We can also add plots to our report:

\begin{Shaded}
\begin{Highlighting}[]
\NormalTok{\textasciigrave{}\textasciigrave{}\textasciigrave{}\{r, out.width = \textquotesingle{}30\%\textquotesingle{}, fig.align = \textquotesingle{}center\textquotesingle{}\}}
\NormalTok{plot(pressure)}
\NormalTok{\textasciigrave{}\textasciigrave{}\textasciigrave{}}
\end{Highlighting}
\end{Shaded}

\begin{center}\includegraphics[width=0.3\linewidth]{Rmarkdown_files/figure-beamer/unnamed-chunk-6-1} \end{center}

The chunk options \texttt{out.width} and \texttt{fig.align} adjust the
size and location
\end{frame}

\begin{frame}[fragile]{R Code Chunks}
\phantomsection\label{r-code-chunks-3}
By default, the code will show up as well. To avoid having the code show
up, you can use an argument. To avoid this, you can use the argument
\texttt{echo=FALSE}. For example:

\begin{Shaded}
\begin{Highlighting}[]
\NormalTok{\textasciigrave{}\textasciigrave{}\textasciigrave{}\{r, echo=FALSE\}}
\NormalTok{summary(pressure)}
\NormalTok{\textasciigrave{}\textasciigrave{}\textasciigrave{}}
\end{Highlighting}
\end{Shaded}

If you want to include the code in the document, but not run the code
and/or include results, you can use the argument \texttt{eval=FALSE}:

\begin{Shaded}
\begin{Highlighting}[]
\NormalTok{\textasciigrave{}\textasciigrave{}\textasciigrave{}\{r, eval=FALSE\}}
\NormalTok{\# You can install the tidyverse using:}
\NormalTok{install.packages(tidyverse)}
\NormalTok{\textasciigrave{}\textasciigrave{}\textasciigrave{}}
\end{Highlighting}
\end{Shaded}
\end{frame}

\begin{frame}[fragile]{R Code Chunks (global \texttt{knitr} options)}
\phantomsection\label{r-code-chunks-global-knitr-options}
One of the R chunks may contain a complex looking call:

\begin{Shaded}
\begin{Highlighting}[]
\NormalTok{\textasciigrave{}\textasciigrave{}\textasciigrave{}\{r setup, include=FALSE\}}
\NormalTok{knitr::opts\_chunk$set(echo = TRUE)}
\NormalTok{\textasciigrave{}\textasciigrave{}\textasciigrave{}}
\end{Highlighting}
\end{Shaded}

The \texttt{include=FALSE} option in the chunk call will tell the system
to run the R code, but not include it in the html/pdf/word report.

The R code in the chunk sets the \texttt{knitr} default to
\texttt{echo=TRUE} in the call for any R chunks following this one.
\end{frame}

\begin{frame}[fragile]{Other Code Chunks}
\phantomsection\label{other-code-chunks}
\texttt{knitr} can execute code in many languages besides R:

\begin{Shaded}
\begin{Highlighting}[]
\NormalTok{\textasciigrave{}\textasciigrave{}\textasciigrave{}\{python\}}
\NormalTok{x = [2, 1, 3]}
\NormalTok{print(x[0],"Hello Python!")}
\NormalTok{\textasciigrave{}\textasciigrave{}\textasciigrave{}}
\end{Highlighting}
\end{Shaded}

\begin{verbatim}
## 2 Hello Python!
\end{verbatim}

Some of the available language engines include: Python, SQL, Bash, Rcpp,
Stan, JavaScript, and CSS.
\end{frame}

\begin{frame}[fragile]{\LaTeX~Equations}
\phantomsection\label{equations}
One exciting feature about R markdown is that it allows you to include
\LaTeX~equations to be rendered in html, pdf, or Word.

For example, to show the probability under the normal curve in the
interval \((a,b)\), you can use the \LaTeX~code: \footnotesize

\begin{Shaded}
\begin{Highlighting}[]
\NormalTok{\textbackslash{}mbox\{Pr\}(a \textless{} x \textless{} b) = }
\NormalTok{  \textbackslash{}int\_a\^{}b \textbackslash{}frac\{1\}\{\textbackslash{}sqrt\{2\textbackslash{}pi\}\textbackslash{}sigma\} }
\NormalTok{  e\^{}\{{-}\textbackslash{}frac\{1\}\{2\}\textbackslash{}left( \textbackslash{}frac\{x{-}\textbackslash{}mu\}\{\textbackslash{}sigma\} \textbackslash{}right)\^{}2\} \textbackslash{}, dx}
\end{Highlighting}
\end{Shaded}

\normalsize

This will be rendered by \texttt{knitr} as:

\[\mbox{Pr}(a < x < b) = \int_a^b \frac{1}{\sqrt{2\pi}\sigma} e^{-\frac{1}{2}\left( \frac{x-\mu}{\sigma} \right)^2} \, dx\]
\end{frame}

\begin{frame}[fragile]{Other Options (tabsets)}
\phantomsection\label{other-options-tabsets}
One very useful tool in Markdown is to add tabs to the html document:

\begin{Shaded}
\begin{Highlighting}[]
\NormalTok{\#\# This is a header \{.tabset\}}
\NormalTok{\#\#\# This is Tab 1}
\NormalTok{\#\#\# This is Tab 2}
\NormalTok{\#\# This is a new header (not in the tabs)}
\end{Highlighting}
\end{Shaded}

The \texttt{\{.tabset\}} option will make all lower level subheadings
tabs in the html. The next same or higher level header will break out of
the tabset.
\end{frame}

\begin{frame}[fragile]{Other Options (tabsets)}
\phantomsection\label{other-options-tabsets-1}
\textbf{Pro tip:} tabs can be generated dynamically:

\begin{Shaded}
\begin{Highlighting}[]
\NormalTok{\#\# My Header \{.tabset\}}
\NormalTok{\textasciigrave{}\textasciigrave{}\textasciigrave{}\{r, results = \textquotesingle{}asis\textquotesingle{}\}}
\NormalTok{n \textless{}{-} 10}
\NormalTok{for (i in 1:n)\{}
\NormalTok{  cat("\#\#\#" , "Tab", i, "\textbackslash{}n")}
\NormalTok{    print(i)}
\NormalTok{  cat("\textbackslash{}n\textbackslash{}n")}
\NormalTok{\}}
\NormalTok{\textasciigrave{}\textasciigrave{}\textasciigrave{}}
\end{Highlighting}
\end{Shaded}
\end{frame}

\begin{frame}{Other Options}
\phantomsection\label{other-options}
We will explore other R markdown options: headers, tabsets, and
\LaTeX~equations in class and in your Exercises.

\underline{Note:} From now on, all R-based Exercises should preferably
be done using an R markdown document. Your document should include
headers, descriptive text, equations, R code, and plots/figures!
\end{frame}

\begin{frame}{More on R markdown}
\phantomsection\label{more-on-r-markdown}
There is a lot more you can do with R markdown. We highly recommend you
continue learning as you gain more experience writing reports in R.
There are many free resources on the internet including:

\begin{enumerate}
\tightlist
\item
  RStudio's tutorial: \url{https://rmarkdown.rstudio.com}
\item
  The cheat sheet:
  \url{https://www.rstudio.com/wp-content/uploads/2015/02/rmarkdown-cheatsheet.pdf}
\item
  The knitR book: \url{https://yihui.name/knitr/}
\end{enumerate}
\end{frame}

\begin{frame}{More Practice}
\phantomsection\label{more-practice}
R markdown is a powerful tool for literate prgramming. To gain more
practice, as part of your homework you will recreate the .Rmd file for
the example file in the homework folder: ``Rmarkdown\_example.html.''
\end{frame}

\begin{frame}[fragile]{Install the following packages:}
\phantomsection\label{install-the-following-packages}
For the next lecture on tidyverse R packages please install the
following:

The tidyverse:

\begin{Shaded}
\begin{Highlighting}[]
\FunctionTok{install.packages}\NormalTok{(}\StringTok{"tidyverse"}\NormalTok{)}
\end{Highlighting}
\end{Shaded}

R Packages:

\begin{Shaded}
\begin{Highlighting}[]
\FunctionTok{install.packages}\NormalTok{(}\FunctionTok{c}\NormalTok{(}\StringTok{"devtools"}\NormalTok{,}\StringTok{"styler"}\NormalTok{,}\StringTok{"testthat"}\NormalTok{))}
\end{Highlighting}
\end{Shaded}

Shiny:

\begin{Shaded}
\begin{Highlighting}[]
\FunctionTok{install.packages}\NormalTok{(}\StringTok{"shiny"}\NormalTok{)}
\end{Highlighting}
\end{Shaded}
\end{frame}

\begin{frame}[fragile]{Session info}
\phantomsection\label{session-info}
\tiny

\begin{Shaded}
\begin{Highlighting}[]
\FunctionTok{sessionInfo}\NormalTok{()}
\end{Highlighting}
\end{Shaded}

\begin{verbatim}
## R version 4.5.2 (2025-10-31 ucrt)
## Platform: x86_64-w64-mingw32/x64
## Running under: Windows 11 x64 (build 22631)
## 
## Matrix products: default
##   LAPACK version 3.12.1
## 
## locale:
## [1] LC_COLLATE=English_United States.utf8 
## [2] LC_CTYPE=English_United States.utf8   
## [3] LC_MONETARY=English_United States.utf8
## [4] LC_NUMERIC=C                          
## [5] LC_TIME=English_United States.utf8    
## 
## time zone: America/New_York
## tzcode source: internal
## 
## attached base packages:
## [1] stats     graphics  grDevices utils     datasets  methods   base     
## 
## loaded via a namespace (and not attached):
##  [1] digest_0.6.39     fastmap_1.2.0     xfun_0.55         Matrix_1.7-4     
##  [5] lattice_0.22-7    reticulate_1.44.1 knitr_1.51        htmltools_0.5.9  
##  [9] png_0.1-8         rmarkdown_2.30    cli_3.6.5         grid_4.5.2       
## [13] withr_3.0.2       compiler_4.5.2    rprojroot_2.1.1   here_1.0.2       
## [17] rstudioapi_0.17.1 tools_4.5.2       evaluate_1.0.5    Rcpp_1.1.1       
## [21] yaml_2.3.12       otel_0.2.0        jsonlite_2.0.0    rlang_1.1.7
\end{verbatim}
\end{frame}

\end{document}
